% Generated draft: requires mathematical and source review.
\documentclass[11pt]{article}
\usepackage[margin=25mm]{geometry}
\usepackage{fontspec}
\setmainfont{Latin Modern Roman}
\usepackage{amsmath,amssymb,mathtools,unicode-math}
\setmathfont{Latin Modern Math}
\usepackage{enumitem,needspace,xurl,hyperref}
\setlength{\parindent}{0pt}
\setlength{\parskip}{5pt}
\setlength{\emergencystretch}{3em}
\newcommand{\N}{\mathbb{N}}
\newcommand{\Z}{\mathbb{Z}}
\newcommand{\Q}{\mathbb{Q}}
\newcommand{\R}{\mathbb{R}}
\newcommand{\C}{\mathbb{C}}
\newcommand{\F}{\mathbb{F}}
\newcommand{\A}{\mathbb{A}}
\newcommand{\PP}{\mathbb{P}}
\newcommand{\E}{\mathbb{E}}
\newcommand{\eps}{\varepsilon}
\newcommand{\dd}{\,\mathrm{d}}
\begin{document}
\section*{UnsolvedMath ID: 30001165}
\textbf{Source problem code:} catalogue entry
\subsection{Definitions and mathematical statement}
Let \(\Omega\subset\R^n\) denote the domain mentioned in the source. The source does not specify the dimension range, regularity or boundedness of \(\Omega\), boundary conditions, potential class, or eigenvalue-moment exponent. Consequently, it does not determine a single quantified conjecture.

To describe the requested estimate conditionally, write \(H_{\Omega,V}=-\Delta+V\) for a self-adjoint Schrödinger realization once its boundary conditions and hypotheses on the real potential \(V\) have been specified. Set \(V_-(x)=\max\{-V(x),0\}\). If the negative spectrum consists of discrete eigenvalues \(\lambda_j(H_{\Omega,V})<0\), counted with multiplicity, define
\[
M_\gamma(H_{\Omega,V})=\sum_{\lambda_j<0}|\lambda_j|^\gamma\quad(\gamma>0),
\qquad
M_0(H_{\Omega,V})=\#\{j:\lambda_j<0\}.
\]
These definitions supply notation; they do not impose hypotheses absent from the source.

\textbf{Source question.} Can Lieb--Thirring estimates for a Schrödinger operator on \(\Omega\) be improved relative to the corresponding estimates on \(\R^n\), for example through a boundary correction? Schematically, in a setting where a whole-space constant \(L_{\gamma,n}\) gives the relevant comparison, the question seeks an improvement of an estimate of the form
\[
M_\gamma(H_{\Omega,V})\le L_{\gamma,n}\int_\Omega V_-(x)^{\gamma+n/2}\,\dd x,
\]
possibly by reducing the constant or by including a nontrivial correction reflecting the boundary. This display describes the intended type of estimate, not an assertion that it holds for unspecified boundary conditions, dimensions, exponents, or potentials. The form, sign, size, and uniformity of a boundary correction, and the meaning of an improved constant, remain unspecified.

The supplied background also discusses a logarithmically corrected two-dimensional Cwikel--Lieb--Rozenblum question and magnetic-operator questions in an extracted composite. Their formulations are absent, so they cannot be incorporated as precise additional mathematical statements here. The supplied dated assessment labels that composite partially solved; it does not supply a resolution of this domain-boundary question.
\subsection{Short English statement}
The question asks whether the geometry or boundary of a domain \(\Omega\) permits stronger Lieb--Thirring bounds for negative Schrödinger eigenvalues than the corresponding whole-space bounds. Boundary conditions, admissible domains and potentials, exponents, and the proposed correction must be specified before this becomes a precise conjecture. The supplied partial-solution label also refers to other questions whose statements are missing.
\subsection{Sources}
\begin{itemize}
\item Local UnsolvedMath registry, supplied record with permanent ID \(30001165\), title \emph{Boundary Corrections to Domain Lieb--Thirring Constants}, exact source problem code \texttt{catalogue entry}. The record attributes the question to the 2009 Oberwolfach workshop report; no individual proposer is supplied.
\item \emph{Low Eigenvalues of Laplace and Schrödinger Operators} (2009), Oberwolfach Reports, pp.~355--428, DOI \texttt{10.4171/owr/2009/06}. Supplied source citation and literature-source link: \url{https://doi.org/10.4171/owr/2009/06}. The background describes Oberwolfach Reports as documenting mathematical workshops and their open-problem discussions. No problem-specific page or passage is supplied.
\item Kovařík--Vugalter--Weidl, logarithmic estimate, supplied literature-source link: \url{https://doi.org/10.1007/s00220-007-0318-z}. The supplied assessment attributes a solution of the logarithmically corrected two-dimensional CLR question to these authors; the precise estimate is not supplied.
\item Literature checked at: \texttt{2026-08-21}. This is supplied dated curation, not independently verified. Its assessment states that the logarithmically corrected two-dimensional CLR question is solved, while the domain-correction and magnetic-operator questions in the extracted composite remain open in the stated generality. The background expressly calls for independent verification of time-sensitive status and cited claims.
\item The registry labels the statement \texttt{corrected\_verified} and reports that its mathematical content was checked against the official Oberwolfach report, with notation, context, and question form normalized without changing the source-level claim. These are supplied verification claims, not verification performed for this draft. The supplied status is \texttt{partially\_solved}.
\end{itemize}
\end{document}
